% -*- mode : latex; coding: utf-8 -*-
\chapter{Conclusion and Future Work} \label{sec:conclusion}

Inadequate memory capacity is one of the major challenge in modern datacenter. If this problem isn't addressed, memory capacity insufficiency will be the bottleneck of future datacenter performance. Simply adding more DIMMs on processor chip is not scalable due to limited number of pins on chip. One way to resolve this problem is memory disaggregation via CXL interconnection. By adopting this method, we can scale memory capacity and compute power independently. In this study, I proposed cycle-accurate and trace-driven CXL simulator CXLmem. I implemented CXL type 3 device using CXL.mem protocol. Results showed that a CXL attached memory degrades IPC up to 40\%, on average 20\% than local DRAM memory. Read latency of CXL attached memory was almost double of read latency of local DRAM memory. If we use CXL attached memory and local DRAM memory together, we can reduce local DRAM access latency by 18\%.

This research showed impact of CXL attached memory on system. In summary, CXL attached memory is slower than local DRAM memory and using CXL attached memory with local DRAM memory reduces local DRAM read latency. This result shows that if we can divide process data wisely between local DRAM memory and CXL attached memory, it's possible to get better performance with larger memory capacity. How to well divide data is one possible direction for future work. The challenge of this research is that it wasn't possible to exactly validate CXLmem over commercial CXL type 3 device which is not yet provided to public. If commercial CXL type 3 device becomes available, we can validate and do more precise and detail further study. CXLmem is focused on CXL type 3 device with CXL.mem protocol. By adding and modifying some codes, CXLmem can be evolved into CXL type 2 device simulator. 